\documentclass[12pt,a4paper]{article}

\usepackage{iftex}
\ifPDFTeX
  \usepackage[T1]{fontenc}
  \usepackage[utf8]{inputenc}
  \usepackage{newtxtext}
  \usepackage{newtxmath}
\else
  \usepackage{fontspec}
  \setmainfont{Times New Roman}
\fi

\usepackage[brazil]{babel}
\usepackage{geometry}
\usepackage{setspace}
\usepackage{indentfirst}
\usepackage{microtype}
\usepackage{amsmath,amssymb}
\usepackage{booktabs}
\usepackage{longtable}
\usepackage{listings}
\usepackage{xcolor}
\usepackage{hyperref}

\geometry{a4paper,left=3cm,top=3cm,right=2cm,bottom=2cm}
\onehalfspacing
\setlength{\parindent}{1.25cm}
\setlength{\parskip}{0pt}
\hypersetup{colorlinks=true,linkcolor=black,citecolor=black,urlcolor=black}

\lstdefinestyle{pythonstyle}{
  language=Python,
  basicstyle=\ttfamily\small,
  keywordstyle=\color{blue!60!black},
  commentstyle=\color{green!40!black},
  stringstyle=\color{red!50!black},
  showstringspaces=false,
  breaklines=true,
  frame=single,
  tabsize=4
}
\lstset{style=pythonstyle}

\newcommand{\aluno}{Nome Completo do Aluno}
\newcommand{\disciplina}{Metodos Quantitativos em Computacao}
\newcommand{\professor}{Julio Cesar Nievola}
\newcommand{\curso}{Bacharelado em Ciencia da Computacao}
\newcommand{\turma}{Turma U}
\newcommand{\instituicao}{Pontificia Universidade Catolica do Parana}
\newcommand{\titulo}{TDE 02 -- Regressao Linear Multipla com a Base Wine Quality}
\newcommand{\descricao}{Aplicacao de regressao linear multipla para previsao da qualidade de vinhos}

\begin{document}

\begin{titlepage}
\begin{center}
\instituicao

Escola Politecnica

\curso

\vfill
{\bfseries\Large \titulo\par}
\vspace{1cm}
{\large \descricao\par}
\vfill
\begin{flushleft}
\textbf{Aluno:} \aluno

\textbf{Disciplina:} \disciplina

\textbf{Professor:} \professor

\textbf{Turma:} \turma

\textbf{Periodo:} 2026/01
\end{flushleft}
\vfill
Curitiba

2026
\end{center}
\end{titlepage}

\newpage
\section*{Resumo}
\addcontentsline{toc}{section}{Resumo}

Este trabalho apresenta a aplicacao de Regressao Linear Multipla na base Wine Quality. A variavel dependente escolhida foi \texttt{quality}, que representa a nota de qualidade do vinho. Como variaveis explicativas, foram selecionadas as duas variaveis independentes com maior correlacao absoluta com a variavel alvo: \texttt{alcohol} e \texttt{density}. O modelo obtido apresentou coeficiente de determinacao $R^2=0.1974$, indicando que essas duas variaveis explicam aproximadamente 19.74\% da variacao observada na qualidade dos vinhos.

\textbf{Palavras-chave:} regressao linear multipla; correlacao; Wine Quality; coeficiente de determinacao.

\newpage
\tableofcontents

\newpage
\section{Base de Dados}

A base Wine Quality contem 6497 observacoes. Para esta atividade, foram consideradas somente as variaveis numericas. A coluna \texttt{wine\_type}, quando presente, e categorica e foi ignorada no ajuste do modelo.

\begin{longtable}{p{4cm}p{10cm}}
\toprule
\textbf{Variavel} & \textbf{Representacao} \\
\midrule
fixed acidity & Acidez fixa do vinho, formada principalmente por acidos nao volateis. \\
volatile acidity & Acidez volatil, associada principalmente ao acido acetico. \\
citric acid & Quantidade de acido citrico, relacionada ao frescor do vinho. \\
residual sugar & Acucar residual que permanece apos a fermentacao. \\
chlorides & Concentracao de cloretos, relacionada ao teor de sais. \\
free sulfur dioxide & Dioxido de enxofre livre, usado para reduzir oxidacao e crescimento microbiano. \\
total sulfur dioxide & Dioxido de enxofre total, incluindo a parte livre e a combinada. \\
density & Densidade do vinho, influenciada por alcool, acucar e compostos dissolvidos. \\
pH & Medida de acidez ou basicidade do vinho. \\
sulphates & Sulfatos presentes no vinho, associados a conservacao e estabilidade. \\
alcohol & Teor alcoolico do vinho. \\
quality & Nota de qualidade atribuida por avaliadores, usada como variavel dependente. \\
wine_type & Tipo do vinho, tinto ou branco. Coluna categorica ignorada nesta TDE. \\
\bottomrule
\end{longtable}

\section{Variavel Dependente e Variaveis Independentes}

A variavel dependente escolhida foi:

\[
Y = quality
\]

As variaveis independentes numericas consideradas inicialmente foram:

\[
fixed acidity, volatile acidity, citric acid, residual sugar, chlorides, free sulfur dioxide, total sulfur dioxide, density, pH, sulphates, alcohol
\]

\section{Correlacao com a Variavel Dependente}

Foi calculado o coeficiente de correlacao de Pearson entre cada variavel independente e \texttt{quality}.

\begin{table}[h]
\centering
\begin{tabular}{lr}
\toprule
\textbf{Variavel independente} & \textbf{Correlacao com quality} \\
\midrule
alcohol & 0.4443 \\
density & -0.3059 \\
volatile acidity & -0.2657 \\
chlorides & -0.2007 \\
citric acid & 0.0855 \\
fixed acidity & -0.0767 \\
free sulfur dioxide & 0.0555 \\
total sulfur dioxide & -0.0414 \\
sulphates & 0.0385 \\
residual sugar & -0.0370 \\
pH & 0.0195 \\
\bottomrule
\end{tabular}
\end{table}

As duas maiores correlacoes em valor absoluto foram obtidas pelas variaveis \texttt{alcohol} e \texttt{density}.

\section{Obtencao da Equacao de Regressao}

A forma geral da Regressao Linear Multipla com duas variaveis explicativas e:

\[
Y = b_0 + b_1X_1 + b_2X_2
\]

Neste trabalho:

\[
X_1 = alcohol
\]

\[
X_2 = density
\]

Logo:

\[
quality = b_0 + b_1 \cdot alcohol + b_2 \cdot density
\]

Os coeficientes estimados pelo metodo dos minimos quadrados foram:

\[
b_0 = 2.8095
\]

\[
b_1 = 0.3246
\]

\[
b_2 = -0.3991
\]

Assim, a equacao final do modelo e:

\[
quality = 2.8095 + 0.3246 \cdot alcohol -0.3991 \cdot density
\]

\section{Avaliacao da Qualidade das Previsoes}

A qualidade das previsoes foi avaliada pelo coeficiente de determinacao, $R^2$:

\[
R^2 = 1 - \frac{SQ_{res}}{SQ_{tot}}
\]

O valor encontrado foi:

\[
R^2 = 0.1974
\]

Portanto, o modelo explica aproximadamente 19.74\% da variacao de \texttt{quality}. Esse resultado indica baixo poder explicativo quando se usam apenas as duas variaveis selecionadas, embora elas sejam as mais correlacionadas individualmente com a qualidade do vinho.

Tambem foram calculadas as seguintes medidas de erro:

\begin{table}[h]
\centering
\begin{tabular}{lr}
\toprule
\textbf{Metrica} & \textbf{Valor} \\
\midrule
MAE & 0.6188 \\
MSE & 0.6119 \\
RMSE & 0.7823 \\
\bottomrule
\end{tabular}
\end{table}

\section{Codigo Fonte}

Os calculos foram realizados pelo arquivo \texttt{tde02\_regressao\_linear\_multipla.py}, presente na raiz do repositorio.

\begin{lstlisting}
python tde02_regressao_linear_multipla.py
\end{lstlisting}

\end{document}
